\begin{enumerate}
	\item \begin{enumerate}
		\item D'après le codage de la figure, le triangle CBF est isocèle en B, et donc les angles à la base ont la même mesure.

		On a donc, en degré:\quad $\widehat{\mathrm{CBF}} = \widehat{\mathrm{BCF}} = 74$.

		\item La somme des angles d'un triangle est un angle plat, de mesure 180\degres.

		On a donc, en degré:\quad $\widehat{\mathrm{CBF}}+\widehat{\mathrm{BCF}}+\widehat{\mathrm{CFB}} = 180$.

		En remplaçant par les mesure connues :\quad $\widehat{\mathrm{CBF}} + 74 + 74 = 180$.

		Finalement, on a donc : \quad $\widehat{\mathrm{CBF}} = 180 - (74 + 74) = 180 - 148 = 32$.

		L'angle $\widehat{\mathrm{CBF}}$ a donc une mesure de 32\degres.
	\end{enumerate}

	\item Dans le triangle ABC, on connaît les longueurs des trois côtés :
	\begin{itemize}
		\item le carré du côté le plus long est :\quad $\mathrm{AC}^2 = 5^2 = 25$;
		\item la somme des carrés des deux petits côtés est :\quad $\mathrm{AB}^2 + \mathrm{BC}^2 = 3^2 + 4^2 = 25$.
	\end{itemize}

	On constate que \quad $\mathrm{AC}^2 = \mathrm{AB}^2 + \mathrm{BC}^2$,\quad d'après la réciproque du théorème de Pythagore, on peut affirmer que le triangle ABC est un triangle rectangle en B.

	\item Le triangle BGF est rectangle en G, donc on peut utiliser les formules de trigonométrie, dans ce triangle.

	On s'intéresse à l'angle $\widehat{\mathrm{FBG}}$.

	\begin{itemize}
		\item On connaît la longueur en cm de l'hypoténuse du triangle :\quad $\mathrm{BF} = 4$;
		\item On connaît la longueur en cm du côté adjacent à l'angle $\widehat{\mathrm{FBG}}$:\quad $\mathrm{BG} = 2$;
	\end{itemize}

	On peut donc calculer le cosinus de l'angle :\quad $\cos\left(\widehat{\mathrm{FBG}}\right) = \dfrac{\text{côté adjacent}}{\text{hypoténuse}} = \dfrac{\mathrm{BG}}{\mathrm{BF}} = \dfrac{2}{4} = \dfrac{1}{2}$.

	On a donc : \quad $\widehat{\mathrm{FBG}} = \arccos\left(\dfrac{1}{2}\right) = 60\degres$.

	\item Les points A, B et G sont alignés si et seulement si l'angle $\widehat{\mathrm{ABG}}$ est un angle plat, dont la mesure est égale à 180\degres.

	Ici on a : \quad $\widehat{\mathrm{ABG}} = \widehat{\mathrm{ABC}} + \widehat{\mathrm{CBF}}+\widehat{\mathrm{FBG}}$.

	Avec les mesures connues en degré : \quad $\widehat{\mathrm{ABG}} = 90 + 32 + 60 = 182$.

	L'angle $\widehat{\mathrm{ABG}}$ n'est pas un angle plat, les points ne sont donc pas alignés.
\end{enumerate}

\needspace{10\baselineskip}
% saut de page si moins de 10 lignes disponibles
